\documentclass[10pt,a4paper]{amsart}
\usepackage{amsmath,amssymb,amsthm,graphicx}
\usepackage{verbatim}
\newenvironment{reference}{\comment}{\endcomment}
%\newenvironment{reference}{}{}
\newenvironment{slogan}{\comment}{\endcomment}
\newenvironment{history}{\comment}{\endcomment}

% For commutative diagrams we use Xy-pic
\usepackage[all]{xy}

% We use 2cell for 2-commutative diagrams.
\xyoption{2cell}
% Theorem environments.
%
\theoremstyle{plain}
\newtheorem{theorem}[subsection]{Theorem}
\newtheorem{proposition}[subsection]{Proposition}
\newtheorem{lemma}[subsection]{Lemma}

\theoremstyle{definition}
\newtheorem{definition}[subsection]{Definition}
\newtheorem{example}[subsection]{Example}
\newtheorem{exercise}[subsection]{Exercise}
\newtheorem{situation}[subsection]{Situation}

\theoremstyle{remark}
\newtheorem{remark}[subsection]{Remark}
\newtheorem{remarks}[subsection]{Remarks}

\numberwithin{equation}{subsection}

% Macros
%
\def\lim{\mathop{\mathrm{lim}}\nolimits}
\def\colim{\mathop{\mathrm{colim}}\nolimits}
\def\Spec{\mathop{\mathrm{Spec}}}
\def\Hom{\mathop{\mathrm{Hom}}\nolimits}
\def\Ext{\mathop{\mathrm{Ext}}\nolimits}
\def\SheafHom{\mathop{\mathcal{H}\!\mathit{om}}\nolimits}
\def\SheafExt{\mathop{\mathcal{E}\!\mathit{xt}}\nolimits}
\def\Sch{\mathit{Sch}}
\def\Mor{\mathop{\mathrm{Mor}}\nolimits}
\def\Ob{\mathop{\mathrm{Ob}}\nolimits}
\def\Sh{\mathop{\mathit{Sh}}\nolimits}
\def\NL{\mathop{N\!L}\nolimits}
\def\CH{\mathop{\mathrm{CH}}\nolimits}
\def\proetale{{pro\text{-}\acute{e}tale}}
\def\etale{{\acute{e}tale}}
\def\QCoh{\mathit{QCoh}}
\def\Ker{\mathop{\mathrm{Ker}}}
\def\Im{\mathop{\mathrm{Im}}}
\def\Coker{\mathop{\mathrm{Coker}}}
\def\Coim{\mathop{\mathrm{Coim}}}

% Boxtimes
%
\DeclareMathSymbol{\boxtimes}{\mathbin}{AMSa}{"02}

%
% Macros for moduli stacks/spaces
%
\def\QCohstack{\mathcal{QC}\!\mathit{oh}}
\def\Cohstack{\mathcal{C}\!\mathit{oh}}
\def\Spacesstack{\mathcal{S}\!\mathit{paces}}
\def\Quotfunctor{\mathrm{Quot}}
\def\Hilbfunctor{\mathrm{Hilb}}
\def\Curvesstack{\mathcal{C}\!\mathit{urves}}
\def\Polarizedstack{\mathcal{P}\!\mathit{olarized}}
\def\Complexesstack{\mathcal{C}\!\mathit{omplexes}}
% \Pic is the operator that assigns to X its picard group, usage \Pic(X)
% \Picardstack_{X/B} denotes the Picard stack of X over B
% \Picardfunctor_{X/B} denotes the Picard functor of X over B
\def\Pic{\mathop{\mathrm{Pic}}\nolimits}
\def\Picardstack{\mathcal{P}\!\mathit{ic}}
\def\Picardfunctor{\mathrm{Pic}}
\def\Deformationcategory{\mathcal{D}\!\mathit{ef}}
\usepackage{cleveref}
\crefname{originallemma}{original lemma}{original lemmas}
\Crefname{originallemma}{Original lemma}{Original lemmas}
\makeatletter
\newlabel{categories-lemma-characterize-mono-epi}{{08LR}{source}{}{sourcepointer-1}{}}
\newlabel{categories-lemma-characterize-mono-epi@cref}{{[originallemma][0][]08LR}{[1][0][]source}{}{}{}}
\makeatother
\crefname{originallemma}{original lemma}{original lemmas}
\Crefname{originallemma}{Original lemma}{Original lemmas}
\makeatletter
\newlabel{lemma-coequalizer-surjection}{{086K}{source}{}{sourcepointer-2}{}}
\newlabel{lemma-coequalizer-surjection@cref}{{[originallemma][0][]086K}{[1][0][]source}{}{}{}}
\makeatother
\crefname{originallemma}{original lemma}{original lemmas}
\Crefname{originallemma}{Original lemma}{Original lemmas}
\makeatletter
\newlabel{categories-lemma-adjoint-fully-faithful}{{07RB}{source}{}{sourcepointer-3}{}}
\newlabel{categories-lemma-adjoint-fully-faithful@cref}{{[originallemma][0][]07RB}{[1][0][]source}{}{}{}}
\makeatother
\title{Implications between exactness properties of a morphism of topoi}
\author{The Stacks Project Authors}
\date{stacks-project:sites.tex:local-snapshot-20261012; GNU FDL 1.2 or later}
\begin{document}
\maketitle
\noindent\textit{Editorial note.} All delivered mathematical text is the author's own, copied from the pinned Stacks Project chapter file sites-cohomology.tex; no statement or proof step was rewritten, reordered or removed.\par\medskip
\noindent\textit{Editorial note.} The proof is the author's complete argument for the eight groups of implications; the sentence 'This is immediate from the definitions' concerns the first group, whose content is a definitional consequence, and the last group is attributed to a named categorical result, both preserved verbatim inside the counted span.\par\medskip
\noindent\textit{Editorial note.} Every labelled reference occurring inside the counted spans is carried in assembly.external\_references as a non-load-bearing navigation pointer (same-chapter ones at their definition site, cross-chapter ones at the invocation inside the counted span), so no load-bearing mathematics is replaced by a stub. No bibliography citation occurs inside the counted spans.\par\medskip
\noindent\textit{Editorial note.} The checked chapter list for same-name collisions is the set of stacks chapters staged in this repository: chow, cohomology, divisors, etale-cohomology, flat, more-algebra, more-morphisms, morphisms, spaces-more-morphisms, varieties together with this card's chapter sites-cohomology.\par\medskip
\noindent\textit{Editorial note.} The preamble is the chapter's own retained minimal preamble (sources/stacks-sites-cohomology/retained-preamble.tex), a verbatim excerpt of the native preamble: the Xy-pic setup and the native theorem declarations with the mathematical macros.\par\medskip
\noindent\textit{Allowed standard prerequisites (editorial).}\begin{itemize}
\item Characterization of monomorphisms and epimorphisms by the diagonal and the codiagonal. Categories, Lemma categories-lemma-characterize-mono-epi. Rewrites injectivity and surjectivity as isomorphisms statements used throughout the proof.
\item Full faithfulness of an adjoint pair. Categories, Lemma categories-lemma-adjoint-fully-faithful. Supplies the equivalence between the counit and unit being an isomorphism and full faithfulness.
\item The pullback of a morphism of topoi is exact. Same chapter, the definition of a morphism of topoi. Allows the pullback to be applied to the diagrams considered in the proof.
\end{itemize}\medskip
\ifdefined\hypertarget\hypertarget{sourcepointer-1}{}\fi
\noindent\textit{Source reference.} Original lemma 08LR: sources/stacks-sites/pinned-original-source.tex, source lines 9867--9867. Navigation pointer invoked at proof line 9867; the counted span contains the pointer only and the cited result is not reproduced inside it, so the pointer is registered as a non-load-bearing navigation stub in the established form for Stacks material. No bibliography citation is involved.\par\medskip
\ifdefined\hypertarget\hypertarget{sourcepointer-2}{}\fi
\noindent\textit{Source reference.} Original lemma 086K: sources/stacks-sites/pinned-original-source.tex, source lines 9883--9883. Navigation pointer invoked at proof line 9883; the counted span contains the pointer only and the cited result is not reproduced inside it, so the pointer is registered as a non-load-bearing navigation stub in the established form for Stacks material. No bibliography citation is involved.\par\medskip
\ifdefined\hypertarget\hypertarget{sourcepointer-3}{}\fi
\noindent\textit{Source reference.} Original lemma 07RB: sources/stacks-sites/pinned-original-source.tex, source lines 9959--9959. Navigation pointer invoked at proof line 9959; the counted span contains the pointer only and the cited result is not reproduced inside it, so the pointer is registered as a non-load-bearing navigation stub in the established form for Stacks material. No bibliography citation is involved.\par\medskip
\begin{lemma}
\label{lemma-exactness-properties}
Let $f : \Sh(\mathcal{C}) \to \Sh(\mathcal{D})$ be
a morphism of topoi. Consider the following properties (on sheaves
of sets):
\begin{enumerate}
\item $f_*$ is faithful,
\item $f_*$ is fully faithful,
\item $f^{-1}f_*\mathcal{F} \to \mathcal{F}$ is surjective for
all $\mathcal{F}$ in $\Sh(\mathcal{C})$,
\item $f_*$ transforms surjections into surjections,
\item $f_*$ commutes with coequalizers,
\item $f_*$ commutes with pushouts,
\item $f^{-1}f_*\mathcal{F} \to \mathcal{F}$ is an isomorphism for
all $\mathcal{F}$ in $\Sh(\mathcal{C})$,
\item $f_*$ reflects injections,
\item $f_*$ reflects surjections,
\item $f_*$ reflects bijections, and
\item for any surjection $\mathcal{F} \to f^{-1}\mathcal{G}$ there
exists a surjection $\mathcal{G}' \to \mathcal{G}$ such that
$f^{-1}\mathcal{G}' \to f^{-1}\mathcal{G}$ factors through
$\mathcal{F} \to f^{-1}\mathcal{G}$.
\end{enumerate}
Then we have the following implications
\begin{enumerate}
\item[(a)] (2) $\Rightarrow$ (1),
\item[(b)] (3) $\Rightarrow$ (1),
\item[(c)] (7) $\Rightarrow$ (1), (2), (3), (8), (9), (10).
\item[(d)] (3) $\Leftrightarrow$ (9),
\item[(e)] (6) $\Rightarrow$ (4) and (5) $\Rightarrow$ (4),
\item[(f)] (4) $\Leftrightarrow$ (11),
\item[(g)] (9) $\Rightarrow$ (8), (10), and
\item[(h)] (2) $\Leftrightarrow$ (7).
\end{enumerate}
Picture
$$
\xymatrix{
(6) \ar@{=>}[rd] & & & & & (9) \ar@{=>}[r] \ar@{=>}[rd] & (8) \\
& (4) \ar@{<=>}[r] & (11) &
(2) \ar@{<=>}[r] &
(7) \ar@{=>}[ru] \ar@{=>}[rd] & & (10) \\
(5) \ar@{=>}[ur] & & & & & (3) \ar@{=>}[r] & (1)
}
$$
\end{lemma}

\begin{proof}
Proof of (a): This is immediate from the definitions.

\medskip\noindent
Proof of (b). Suppose that $a, b : \mathcal{F} \to \mathcal{F}'$ are
maps of sheaves on $\mathcal{C}$. If $f_*a = f_*b$, then
$f^{-1}f_*a = f^{-1}f_*b$. Consider the commutative diagram
$$
\xymatrix{
\mathcal{F} \ar@<-1ex>[r] \ar@<1ex>[r] & \mathcal{F}' \\
f^{-1}f_*\mathcal{F} \ar@<-1ex>[r] \ar@<1ex>[r] \ar[u] &
f^{-1}f_*\mathcal{F}' \ar[u]
}
$$
If the bottom two arrows are equal and the vertical arrows are surjective
then the top two arrows are equal. Hence (b) follows.

\medskip\noindent
Proof of (c). Suppose that $a : \mathcal{F} \to \mathcal{F}'$ is a
map of sheaves on $\mathcal{C}$. Consider the commutative diagram
$$
\xymatrix{
\mathcal{F} \ar[r] & \mathcal{F}' \\
f^{-1}f_*\mathcal{F} \ar[r] \ar[u] &
f^{-1}f_*\mathcal{F}' \ar[u]
}
$$
If (7) holds, then the vertical arrows are isomorphisms.
Hence if $f_*a$ is injective (resp.\ surjective, resp.\ bijective)
then the bottom arrow is injective (resp.\ surjective, resp.\ bijective) and
hence the top arrow is injective (resp.\ surjective, resp.\ bijective).
Thus we see that (7) implies (8), (9), (10). It is clear that (7) implies (3).
The implications (7) $\Rightarrow$ (2), (1) follow from (a) and (h) which
we will see below.

\medskip\noindent
Proof of (d). Assume (3). Suppose that $a : \mathcal{F} \to \mathcal{F}'$
is a map of sheaves on $\mathcal{C}$ such that $f_*a$ is surjective.
As $f^{-1}$ is exact this implies that
$f^{-1}f_*a : f^{-1}f_*\mathcal{F} \to f^{-1}f_*\mathcal{F}'$
is surjective. Combined with (3) this implies that $a$ is surjective.
This means that (9) holds.
Assume (9). Let $\mathcal{F}$ be a sheaf on $\mathcal{C}$.
We have to show that the map $f^{-1}f_*\mathcal{F} \to \mathcal{F}$ is
surjective. It suffices to show that
$f_*f^{-1}f_*\mathcal{F} \to f_*\mathcal{F}$ is surjective.
And this is true because there is a canonical map
$f_*\mathcal{F} \to f_*f^{-1}f_*\mathcal{F}$ which is a one-sided inverse.

\medskip\noindent
Proof of (e). We use
Categories, Lemma \ref{categories-lemma-characterize-mono-epi}
without further mention.
If $\mathcal{F} \to \mathcal{F}'$ is surjective then
$\mathcal{F}' \amalg_\mathcal{F} \mathcal{F}' \to \mathcal{F}'$
is an isomorphism. Hence (6) implies that
$$
f_*\mathcal{F}' \amalg_{f_*\mathcal{F}} f_*\mathcal{F}' =
f_*(\mathcal{F}' \amalg_\mathcal{F} \mathcal{F}')
\longrightarrow
f_*\mathcal{F}'
$$
is an isomorphism also. And this in turn implies that
$f_*\mathcal{F} \to f_*\mathcal{F}'$ is surjective.
Hence we see that (6) implies (4). If $\mathcal{F} \to \mathcal{F}'$
is surjective then $\mathcal{F}'$ is the coequalizer of the two
projections $\mathcal{F} \times_{\mathcal{F}'} \mathcal{F} \to \mathcal{F}$
by Lemma \ref{lemma-coequalizer-surjection}.
Hence if (5) holds, then $f_*\mathcal{F}'$ is the coequalizer of
the two projections
$$
f_*(\mathcal{F} \times_{\mathcal{F}'} \mathcal{F}) =
f_*\mathcal{F} \times_{f_*\mathcal{F}'} f_*\mathcal{F}
\longrightarrow
f_*\mathcal{F}
$$
which clearly means that $f_*\mathcal{F} \to f_*\mathcal{F}'$ is
surjective. Hence (5) implies (4) as well.

\medskip\noindent
Proof of (f). Assume (4). Let $\mathcal{F} \to f^{-1}\mathcal{G}$ be
a surjective map of sheaves on $\mathcal{C}$. By (4) we see that
$f_*\mathcal{F} \to f_*f^{-1}\mathcal{G}$ is surjective. Let
$\mathcal{G}'$ be the fibre product
$$
\xymatrix{
f_*\mathcal{F} \ar[r] & f_*f^{-1}\mathcal{G} \\
\mathcal{G}' \ar[u] \ar[r] & \mathcal{G} \ar[u]
}
$$
so that $\mathcal{G}' \to \mathcal{G}$ is surjective also. Consider
the commutative diagram
$$
\xymatrix{
\mathcal{F} \ar[r] & f^{-1}\mathcal{G} \\
f^{-1}f_*\mathcal{F} \ar[r] \ar[u] & f^{-1}f_*f^{-1}\mathcal{G} \ar[u] \\
f^{-1}\mathcal{G}' \ar[u] \ar[r] & f^{-1}\mathcal{G} \ar[u]
}
$$
and we see the required result. Conversely, assume (11).
Let $a : \mathcal{F} \to \mathcal{F}'$ be surjective map of sheaves
on $\mathcal{C}$. Consider the fibre product diagram
$$
\xymatrix{
\mathcal{F} \ar[r] & \mathcal{F}' \\
\mathcal{F}'' \ar[u] \ar[r] & f^{-1}f_*\mathcal{F}' \ar[u]
}
$$
Because the lower horizontal arrow is surjective and by (11) we can
find a surjection $\gamma : \mathcal{G}' \to f_*\mathcal{F}'$ such
that $f^{-1}\gamma$ factors through $\mathcal{F}'' \to f^{-1}f_*\mathcal{F}'$:
$$
\xymatrix{
& \mathcal{F} \ar[r] & \mathcal{F}' \\
f^{-1}\mathcal{G}' \ar[r] & \mathcal{F}'' \ar[u] \ar[r] &
f^{-1}f_*\mathcal{F}' \ar[u]
}
$$
Pushing this down using $f_*$ we get a commutative diagram
$$
\xymatrix{
& f_*\mathcal{F} \ar[r] & f_*\mathcal{F}' \\
f_*f^{-1}\mathcal{G}' \ar[r] & f_*\mathcal{F}'' \ar[u] \ar[r] &
f_*f^{-1}f_*\mathcal{F}' \ar[u] \\
\mathcal{G}' \ar[u] \ar[rr] & & f_*\mathcal{F}' \ar[u]
}
$$
which proves that (4) holds.

\medskip\noindent
Proof of (g). Assume (9). We use
Categories, Lemma \ref{categories-lemma-characterize-mono-epi}
without further mention. Let $a : \mathcal{F} \to \mathcal{F}'$ be a map
of sheaves on $\mathcal{C}$ such that $f_*a$ is injective. This means that
$f_*\mathcal{F} \to f_*\mathcal{F} \times_{f_*\mathcal{F}'} f_*\mathcal{F} =
f_*(\mathcal{F} \times_{\mathcal{F}'} \mathcal{F})$ is an isomorphism.
Thus by (9) we see that
$\mathcal{F} \to \mathcal{F} \times_{\mathcal{F}'} \mathcal{F}$ is
surjective, i.e., an isomorphism. Thus $a$ is injective, i.e., (8) holds.
Since (10) is trivially equivalent to (8) $+$ (9) we are done with (g).

\medskip\noindent
Proof of (h). This is
Categories, Lemma \ref{categories-lemma-adjoint-fully-faithful}.
\end{proof}

\end{document}
